\documentclass[11pt]{article}
\usepackage[margin=1in]{geometry}
\usepackage{times}
% \usepackage[skip=0pt]{parskip}
\pagestyle{empty}

\begin{document}

\noindent To the Editors\\
International Journal of Molecular Sciences

\vspace{1.5em}
\noindent Dear Editor,

We are pleased to submit our manuscript entitled ``A Framework for Global Background Subtraction in Fiber Diffraction Patterns from Striated Muscle'' for consideration for publication in the International Journal of Molecular Sciences, possibly in the Topic Collection  ``Feature Papers in Molecular Biophysics''.

X-ray fiber diffraction is a technique that provides unique insight into the sarcomere-level molecular structure of muscle and has proven valuable for studying the structural consequences of cardiac and skeletal myopathies. Analysis of muscle diffraction patterns can therefore contribute to understanding the molecular mechanisms underlying diseases ranging from nemaline myopathy to heart failure. The quantitative interpretation of these diffraction patterns, however, depends on the reliable separation of weak diffraction signals from the ordered components of muscle from the strong, spatially non-uniform diffuse background.  While many approaches to accomplishing this exist, there has been no standardized, quantitative basis for choosing the most appropriate two-dimensional background-subtraction methods and optimizing the parameters which introduces subjectivity, risks the loss of faint but meaningful features and undermines reproducibility. 


To address this gap, we present a unified framework that replaces subjective visual assessment with physically-grounded, quantitative evaluation criteria given the lack of ground-truth background data allowing a direct and objective quantitative comparison of the commonly used subtraction methods within a single consistent pipeline.  We show that when applied to synchrotron diffraction patterns from striated muscle, the optimized parameters substantially improve the consistency of background removal and reduce artifacts relative to default or ad hoc settings, while preserving key structural features.  We also describe a new iterative procedure to estimate the background from the entire 2D pattern, something not possible with existing approaches.  


This new framework provides both an immediate practical tool and a foundation for ongoing methodological development in quantitative analysis of fiber diffraction patterns from biologicals structures. While primarily developed for muscle diffraction patterns, these tools could be applied to fiber diffraction of other biological systems. The framework has been incorporated into MuscleX, an open-source software package for analysis of muscle fiber diffraction patterns, making it available to the community. This manuscript is original, has not been published previously, and is not under consideration elsewhere. All authors have approved the submission and declare no conflict of interest. We believe the work will be of interest to the journal’s readership in structural and molecular biology, and we thank you for considering it for publication in IJMS.



\vspace{1.0em}
\noindent Sincerely,

\vspace{1.0em}
\noindent Irina Klein, Gady Agam, and Thomas C. Irving\\
Illinois Institute of Technology, Chicago, IL, USA\\
Corresponding authors: irving@illinoistech.edu; iklein@illinoistech.edu

\end{document}
